\section{Maximal does not mean maximum}\label{ch08:sec:maximal}
In the second Ada--Ben example, the greedy procedure produced the matching $\{\text{Ada--}1\}$ and then got stuck, although the matching $\{\text{Ada--}2,\ \text{Ben--}1\}$ solves the problem. Two words describe the difference between these matchings.

\par\noindent\begin{minipage}{\linewidth}
\begin{definition}[Maximal and maximum matchings]\label{ch08:def:maximal}
A matching $M$ is \emph{maximal} if no edge can be added to it: for every edge $e$ of the graph that is not in $M$, the set $M\cup\{e\}$ is not a matching. A matching $M$ is \emph{maximum} if no matching in the graph has more edges than $M$.
\end{definition}
\end{minipage}\par

The matching $\{\text{Ada--}1\}$ is maximal. The graph has two other edges, and each of them shares an endpoint with Ada--$1$: the edge Ada--$2$ shares the endpoint Ada, and the edge Ben--$1$ shares the endpoint~$1$. Adding Ada--$2$ would give Ada two tasks, and adding Ben--$1$ would give task~$1$ two owners. But $\{\text{Ada--}1\}$ is not maximum, because $\{\text{Ada--}2,\ \text{Ben--}1\}$ is a matching with two edges. To get from the first matching to the second, we must remove the edge Ada--$1$ as well as add new edges. A greedy procedure never removes anything, so it can stop at a maximal matching that is not maximum.

Every maximum matching is maximal. Indeed, if some edge $e$ could be added to a matching $M$, then $M\cup\{e\}$ would be a matching with one more edge than $M$, so $M$ would not be maximum. The example shows that the converse is false.

\pauseq{Is the empty matching maximal?}{Only when the graph has no edges. A set consisting of a single edge is always a matching, because there is no second edge with which it could share an endpoint. So if the graph has an edge $e$, then $e$ can be added to the empty matching, which is therefore not maximal. If the graph has no edges, there is nothing to add, and the empty matching is both maximal and maximum.}

\subsection*{Comparison with Chapter~\ref{ch:structure}}
Section~\ref{ch06:sec:order} defined maximal and maximum elements for any partial order. The matchings of a graph are sets of edges, so we can order them by inclusion,~$\subseteq$. In that order, a matching $M$ is maximal when no matching strictly contains it. This agrees with Definition~\ref{ch08:def:maximal}. Suppose first that a matching $M'$ strictly contains $M$, and choose an edge $e$ of $M'$ that is not in $M$. The set $M\cup\{e\}$ lies inside $M'$. Removing edges from a matching cannot create a shared endpoint, so every set of edges contained in a matching is again a matching. Hence $M\cup\{e\}$ is a matching, and $e$ can be added to $M$. Conversely, if some edge $e$ can be added to $M$, then $M\cup\{e\}$ is a matching that strictly contains $M$.

The word \emph{maximum} is a different matter. In Section~\ref{ch06:sec:order}, a maximum lies above every element of the collection, so a maximum matching in that sense would contain every other matching. The Ada--Ben graph has no such matching: a matching containing both $\{\text{Ada--}1\}$ and $\{\text{Ben--}1\}$ would contain two edges with the endpoint~$1$. For matchings, \emph{maximum} always means \emph{largest number of edges}. Two maximum matchings can be quite different. If students $a$ and $b$ both accept tasks $1$ and $2$, then $\{a\text{--}1,\ b\text{--}2\}$ and $\{a\text{--}2,\ b\text{--}1\}$ are both maximum matchings, and neither contains the other. So when an extremal argument (Section~\ref{ch07:sec:walks}) begins ``take a largest matching,'' it must say which comparison it means: largest under inclusion, or largest in size. The first kind can be reached by adding edges one at a time. Reaching the second kind may require undoing earlier choices.

\subsection*{A matching as a function}
A matching also tells each matched student which task to do. For the matching $\{\text{Ada--}1\}$, Ada has task~$1$ and Ben has none, so this is only a partial assignment. When a matching saturates $L$, it gives every student a task, and then it describes a function from $L$ to $R$ in the sense of Chapter~\ref{ch:maps}. More precisely, matchings that saturate $L$ correspond exactly to injective functions $f:L\to R$ such that $s$--$f(s)$ is an edge for every student~$s$.
\begin{itemize}
\item Let $M$ be a matching that saturates $L$. Every student $s$ is matched, so it has exactly one partner; let $f(s)$ be that partner. Then $s$--$f(s)$ is an edge, because it belongs to $M$. The function $f$ is injective: if two different students $s$ and $s'$ had $f(s)=f(s')=t$, then the two different edges $s$--$t$ and $s'$--$t$ of $M$ would share the endpoint $t$.
\item Conversely, let $f:L\to R$ be injective with $s$--$f(s)$ an edge for every $s$. The edges $s$--$f(s)$ for $s\in L$ form a matching that saturates $L$. Indeed, take two of these edges, $s$--$f(s)$ and $s'$--$f(s')$ with $s\ne s'$. A student in $L$ cannot equal a task in $R$, because $L$ and $R$ are disjoint, so a shared endpoint would have to be a shared student, $s=s'$, or a shared task, $f(s)=f(s')$. The first is excluded by the choice of $s$ and $s'$, and the second by injectivity.
\end{itemize}
The two conditions on $f$ do different jobs. Injectivity (Section~\ref{ch03:sec:injsurj}) keeps two students from receiving the same task. The edge condition makes sure that each student receives a task they accept. A function that satisfies only one of the two conditions does not solve the problem.

Keep two sets of edges apart. The edge set of the graph lists the \emph{allowed} pairs; it is fixed by the problem. A matching lists the pairs \emph{currently chosen}; it is a subset of the edge set, and in the next sections it will change as we repair it. Changing the matching never changes which tasks a student accepts.
